\chapter{Chest Pain}

\section*{Definition}
Chest pain is discomfort, pressure, squeezing, burning, sharpness, or heaviness in the anterior chest wall, potentially originating from cardiac, pulmonary, gastrointestinal, musculoskeletal, or psychological structures, ranging from benign to life-threatening etiologies requiring prompt assessment.

\section*{What Happens in Your Body}
Cardiac ischemia activates myocardial sympathetic afferent fibers releasing substance P and CGRP producing retrograde somatic convergence in spinal cord (T1-T4) interpreted as chest pain pericarditis involves parietal pericardium innervated by phrenic nerves producing sharp positional pain pleuritic chest pain arises from parietal pleura innervated by intercostal nerves aggravated by respiration Esophageal pain follows vagal/sympathetic pathways mimicking angina Musculoskeletal pain derives from costosternal junctions, ribs, muscles with reproducible tenderness Anxiety/panic disorders produce chest tightness through hyperventilation and heightened thoracic awareness

\section*{Causes}
\begin{itemize}
  \item Cardiac: Acute coronary syndrome/unstable angina/myocardial infarction, pericarditis, myocarditis, aortic dissection, hypertrophic cardiomyopathy, mitral valve prolapse
  \item Pulmonary: Pulmonary embolism, pneumothorax, pneumonia, pleurisy, COPD exacerbation, pulmonary hypertension
  \item Gastroesophageal: GERD, esophageal spasm, peptic ulcer disease, pancreatitis, cholecystitis
  \item Musculoskeletal: Costochondritis/Tietze syndrome, muscle strain, rib fracture, fibromyalgia
  \item Psychiatric: Panic attack, anxiety, somatization
  \item Other: Herpes zoster (prodromal phase), anxiety-induced chest tightness, factitious pain
\end{itemize}

\section*{When to See a Doctor}
See your doctor if:
\begin{itemize}
  \item Chest pain with diaphoresis, nausea, vomiting, syncope suggesting ACS
  \item Chest pain with sudden onset/dyspnea/hemoptysis suggesting PE
  \item Sudden ``tearing'' chest pain suggesting dissection
  \item Fever/respiratory symptoms suggesting pneumonia
  \item Positional pain relieved/sit forward suggesting pericarditis
  \item Episodic substernal pressure with exertion suggesting angina
  \item New onset in people over 40 years with CV risk factors
\end{itemize}

\section*{Related Symptoms}
\begin{itemize}
  \item Dyspnea
  \item Orthopnea
  \item Paroxysmal nocturnal dyspnea
  \item Palpitations
  \item Dizziness
  \item Presyncope
  \item Syncope
  \item Cough
  \item Hemoptysis
  \item Heartburn
  \item Regurgitation
  \item Odynophagia
  \item Arm/jaw/back radiation
  \item Sweating
  \item Nausea/vomiting
  \item Anxiety
  \item Fear of impending doom
\end{itemize}
\textbf{Important:} Immediate assessment prioritizes ruling out life-threatening causes (ACS, PE, dissection, tension pneumothorax, cardiac tamponade). ECG and serial troponin testing are commonly used when ACS is suspected; D-dimer and chest imaging are selected according to the clinical probability of PE or another diagnosis. Never dismiss chest pain without adequate evaluation regardless of an apparently benign explanation.

\section*{Self-Care / Home Management}
\begin{itemize}
  \item New, sudden, severe, exertional, or unexplained chest pain should not be managed at home; call local emergency services.
  \item Only after a clinician has excluded urgent causes may measures such as rest, ice for a clearly minor chest-wall injury, or an appropriate over-the-counter analgesic be considered.
  \item Do not use self-care advice to postpone assessment of recurrent or unexplained chest pain.
\end{itemize}

\section*{How Your Doctor Will Evaluate You}
\begin{itemize}
  \item Evaluation begins with vital signs, oxygen saturation, focused examination, and an ECG; tests are guided by the suspected diagnosis.
  \item Serial cardiac biomarkers, chest radiography, echocardiography, CT, or stress or coronary testing may be used when clinically indicated.
  \item Risk scores and clinical probability tools can support, but do not replace, clinician assessment.
\end{itemize}

\section*{Treatment Options}
\begin{itemize}
  \item Treatment is diagnosis-specific: emergency reperfusion or antithrombotic treatment for selected coronary syndromes, anticoagulation for selected pulmonary emboli, and urgent surgical or procedural care for dissection or tamponade.
  \item Musculoskeletal, reflux-related, pulmonary, and panic-related pain should be treated only after dangerous causes have been considered and appropriately assessed.
\end{itemize}

\section*{References}
\begin{thebibliography}{99}
\bibitem{chest_pain:ref1} Amsterdam EA, Wenger NK, Brindis RG, et al. ``2014 AHA/ACC guideline for the management of patients with non-ST-elevation acute coronary syndromes.'' J Am Coll Cardiol. 2014;64(24):e139-e228.
\bibitem{chest_pain:ref2} Fihn SD, Gardin JM, Abrams J, et al. ``2012 ACCF/AHA/ACP/AATS/PCNA/SCAI/STS guideline for the diagnosis and management of patients with stable ischemic heart disease.'' Circulation. 2012;126(25):e354-e471.
\bibitem{chest_pain:ref3} Cayley WE Jr. ``Diagnosing the cause of chest pain.'' Am Fam Physician. 2005;72(10):2012-2021.
\bibitem{chest_pain:ref4} Wertli MM, Ruchti KB, Steurer J, et al. ``Diagnostic indicators of non-cardiovascular chest pain: a systematic review and meta-analysis.'' BMC Med. 2013;11:239.
\bibitem{chest_pain:ref5} Libby P, Braunwald E. ``Braunwald's Heart Disease: A Textbook of Cardiovascular Medicine.'' 12th ed. Elsevier; 2022.
\bibitem{chest_pain:ref6} Fuster V, Alexander RW, O'Rourke RA, et al. ``Approach to the patient with chest pain.'' UpToDate. Wolters Kluwer; 2025.
\end{thebibliography}
